% Related work — ~0.5 page main text
% Three buckets: (1) CP theory, (2) heteroscedastic UQ, (3) VEP + conformal

\paragraph{Conformal prediction and conditional coverage.}
Split conformal prediction \citep{vovk2005algorithmic} provides distribution-free marginal coverage, and Mondrian conformal \citep{vovk2003mondrian} extends this to group-conditional coverage via partitioned calibration. \citet{barber2020limits} showed that pointwise conditional coverage $P(Y \in \pset(X) \mid X = x)$ is impossible in finite samples without structural assumptions, motivating coarsened conditioning. CQR \citep{romano2019cqr} uses quantile regression to adapt interval width to local heteroscedasticity in regression; our score Eq.~\eqref{eq:score} is the binary-classification analog. \citet{barber2023exchangeability} generalized CP beyond exchangeability via weighted residual swaps; we use their Theorem~2 for robust forms T1$'$--T3$'$ and T4.

Two recent lines of work target conditional coverage more directly.
\textbf{RLCP} \citep{hore2025rlcp} uses random localization with kernel weights to achieve local coverage at rate $O(n^{-2/(d+2)})$; this degrades in high dimensions ($d > 100$), making it impractical for our genomic feature space ($d \approx 7{,}700$).
\citet{dewolf2025conditional} analyze the conditional validity of heteroscedastic Mondrian conformal for regression, establishing when normalized and Mondrian CP adapt correctly to the noise level; our T3--T5 complement their analysis by deriving the optimal Mondrian granularity (bin count $K$) for classification with a \emph{learned} variance head, with a dimension-free $O(n^{-1/2})$ rate.
Self-Calibrating CP \citep{vanderlaan2024sccp} provides prediction-conditional validity via Venn--Abers calibration, conditioning on the predicted value rather than uncertainty; this is a complementary axis to our $\shat$-bin conditioning.

\paragraph{Heteroscedastic uncertainty quantification.}
Gaussian NLL training for learned variance originates with \citet{nix1994nll}. Deep ensembles \citep{lakshminarayanan2017simple} capture epistemic uncertainty via inter-model disagreement. DEGU \citep{zhou2026degu} combines both: an ensemble of CNN students distilled from a genomic teacher, with Gaussian uncertainty. DEGU's abstract mentions conformal prediction, but the implementation is marginal split CP only --- no class-conditional or Mondrian stratification, no formal local coverage guarantee, and no classification task. Our DEGU-lite ablation (§\ref{sec:experiments}) shows that ensemble $\shat$ provides comparable local coverage on Mendelian (monogenic, strong signal) but $4.3\times$ worse coverage on Complex (polygenic), where the aleatoric component dominates.

\paragraph{Variant effect prediction and benchmarks.}
TraitGym \citep{benegas2025traitgym} standardizes VEP evaluation on Mendelian and complex trait variants, with a chrom-LOO protocol. Input features include CADD \citep{kircher2014cadd}, Borzoi \citep{linder2025borzoi}, and GPN-MSA \citep{benegas2023gpnmsa}. AlphaGenome \citep{avsec2026alphagenome} is the most recent foundation model for genomics, though it has not been evaluated on TraitGym to date. Conformal prediction has been explored in genomic medicine for drug response \citep{papangelou2025cpgenomics} and polygenic risk scores \citep{wang2025cqrprs}, but not for binary variant pathogenicity classification. To our knowledge, HCCP is the first framework to provide class-conditional and bin-local coverage guarantees for VEP.
